\documentclass{tmsce}

% ---------------- Packages ----------------
\usepackage{amsmath,amsfonts,amssymb}
\usepackage{graphicx}
\usepackage{booktabs}
\usepackage{array}
\usepackage{caption}
\usepackage{url}
\usepackage{subfig}
\usepackage{multirow}
\usepackage{xcolor}
\usepackage{enumitem}
\usepackage{longtable}
\usepackage{tikz}
\usetikzlibrary{shapes.geometric, arrows.meta, positioning, calc, fit, backgrounds}

\hyphenation{op-tical net-works semi-conduc-tor trans-crip-tome de-con-vo-lu-tion}

% ---------------- Front-matter  ----------------
\doi{}
\copyrightline{}
\permissions{}
\vol{}
\issue{}
\yearofpub{2026}
\pagerange{}

% ---------------- Title/Author ----------------
\title{Age- and Arm-Stratified Co-expression Network Rewiring in the Murine Kidney During and After Spaceflight: A Comparative Pipeline Analysis}

\authors{Ibrahim M.~Shahid}
\affiliation{Middle College at UNCG \textbullet{} Blue Marble Space Institute of Science \textbullet{} NASA Human Research Program}

\keywords{spaceflight; kidney; transcriptomics; co-expression networks; network rewiring; aging; retinoic acid; LIONESS; node2vec}

\begin{document}
% Remove journal-style footer/header for a preprint-style PDF.
\renewcommand{\journalname}{}
\makeatletter
\renewcommand{\tmsce@topstrip}{}
\makeatother
\fancyfoot[C]{}
\renewcommand{\footrulewidth}{0pt}
\maketitle

\begin{abstract}
Astronauts exhibit elevated nephrolithiasis risk following spaceflight, and prior work links microgravity to distal nephron remodeling and renal transporter dysregulation. Whether spaceflight alters the regulatory coordination among genes---shifting co-expression architecture without necessarily changing mean expression levels---and whether such rewiring differs by age or by spaceflight habitat, remains uncharacterized. We addressed this using OSD-771 (RRRM-2), bulk kidney RNA-seq from 80 female C57BL/6NTac mice under a full $2\times4\times2$ factorial design crossing Age (young, 16 weeks; old, 34 weeks), Environment Group (FLT, GC, VIV, BSL), and Arm (ISS Terminal sacrifice, ISS-T; Live Animal Return, LAR). Nephron segment proportions were estimated via MuSiC deconvolution against a murine single-cell kidney reference and used to remove compositional confounding. Sample-specific co-expression networks were inferred via LIONESS on a shared sparse skeleton, yielding model-predicted contrast-specific networks. Node2vec embeddings, aligned by orthogonal Procrustes across multi-seed runs, produced per-gene rewiring scores for four primary flight contrasts. Focused permutation testing identified 11 significantly rewired genes in young ISS-T animals (FDR $< 0.05$), including Slc6a20a, Cry2, and Gc. Pathway enrichment revealed coordinated ECM and complement remodeling in young ISS-T mice, an EMT and steroid biosynthesis signature in old ISS-T mice, and retinoic acid metabolism as the most statistically robust signal across both arms---peaking in the LAR recovery condition (OR = 19.9, $q = 1.6\times 10^{-5}$). Predictive validation via supervised classification of network-derived topological features achieved above-chance accuracy (Random Forest: 40.0\%, AUC = 0.525), confirming that co-expression topology encodes non-linear spaceflight biology invisible to standard transcriptomic analyses. This report presents a comparative analysis of two pipeline iterations (February and April 2026), documenting methodological advances and the convergent biological picture they reveal.
\end{abstract}

\printkeywords

\tmsceendfrontmatter

% =========================================================
\section{Introduction}

Astronauts experience a roughly five-fold elevated risk of kidney stones (nephrolithiasis) post-flight, with distal convoluted tubule (DCT) ion-transport dysregulation strongly implicated~\cite{whitson2009}. Standard differential expression (DE) tests only detect \textit{mean shifts} in gene expression. Regulatory \textit{rewiring}---changes in how genes co-express and coordinate with one another---is missed by DE when average expression remains stable. Such ``silent'' regulatory restructuring may represent the predominant mode of spaceflight-induced transcriptomic adaptation, particularly in organs like the kidney that must maintain tight homeostatic control over electrolyte and acid-base balance despite radically altered hemodynamic conditions in microgravity.

To address this gap, we developed a multi-phase pipeline integrating cell-type deconvolution, sample-specific network inference, graph embedding, and multi-scale statistical testing. We applied this framework to RRRM-2 (OSD-771), the largest full-factorial kidney RNA-seq spaceflight dataset ($n=80$ mice), and iterated the pipeline across two analytical runs to assess reproducibility, sensitivity gains from methodological refinements, and convergence of biological findings.

The central questions are: (1) whether spaceflight restructures the kidney co-expression network in an age- and arm-dependent manner, (2) which topological hubs and pathways drive that restructuring, and (3) whether these network-derived signals are identifiable and reproducible across pipeline iterations and orthogonal to standard differential expression.


% =========================================================
\section{Experimental Design and Dataset}

\subsection{Full Factorial Design}

The OSD-771 dataset comprises 80 female C57BL/6NTac mice arranged in a $2 \times 4 \times 2$ factorial design:

\begin{itemize}[nosep]
    \item \textbf{Age}: Young (16 weeks at sacrifice) vs.\ Old (34 weeks)
    \item \textbf{Environment Group}: FLT (spaceflight), GC (hardware ground control), VIV (vivarium), BSL (basal/pre-flight)
    \item \textbf{Arm}: ISS-T (terminal sacrifice aboard ISS) vs.\ LAR (live animal return)
\end{itemize}

Each experimental cell (Age $\times$ Arm $\times$ Group) contains 5 mice. The four primary contrasts of interest are:

\begin{enumerate}[nosep]
    \item ISS-T Young: FLT $-$ GC (acute flight effect in young animals)
    \item ISS-T Old: FLT $-$ GC (acute flight effect in old animals)
    \item LAR Young: FLT $-$ GC (post-recovery flight effect in young animals)
    \item LAR Old: FLT $-$ GC (post-recovery flight effect in old animals)
\end{enumerate}


% =========================================================
\section{Analytical Pipeline}

The pipeline integrates advanced network medicine paradigms to resolve context-specific regulatory architectures in small-sample transcriptomics (Figure~\ref{fig:pipeline}).

% ---- Pre-define colors for pipeline diagram ----
\definecolor{prepfill}{RGB}{200,220,245}
\definecolor{prepbord}{RGB}{80,120,180}
\definecolor{preplight}{RGB}{210,228,248}
\definecolor{netfill}{RGB}{180,225,220}
\definecolor{netbord}{RGB}{50,130,120}
\definecolor{netalt}{RGB}{255,220,180}
\definecolor{netaltbord}{RGB}{200,140,50}
\definecolor{inffill}{RGB}{255,210,210}
\definecolor{infbord}{RGB}{180,80,80}
\definecolor{enrfill}{RGB}{195,235,210}
\definecolor{enrbord}{RGB}{60,140,100}
\definecolor{valfill}{RGB}{220,210,240}
\definecolor{valbord}{RGB}{120,90,160}
\definecolor{outfill}{RGB}{230,230,230}
\definecolor{outbord}{RGB}{130,130,130}
\definecolor{tiertext}{RGB}{100,100,120}
\definecolor{arrcolor}{RGB}{100,100,110}
\definecolor{grpblue}{RGB}{220,230,245}
\definecolor{grpteal}{RGB}{220,240,235}
\definecolor{grprose}{RGB}{245,225,225}

% ---- TikZ pipeline flowchart ----
\begin{figure*}[!ht]
\centering
\resizebox{\textwidth}{!}{%
\begin{tikzpicture}[
    % -- base node style --
    pbox/.style={
        rectangle, rounded corners=4pt,
        minimum width=2.8cm, minimum height=1.5cm,
        text width=2.6cm, align=center, font=\small\sffamily,
        line width=0.8pt,
    },
    % tier labels
    tierlabel/.style={font=\footnotesize\sffamily\itshape, text=tiertext},
    % arrows
    arr/.style={-{Stealth[length=5pt,width=4pt]}, line width=0.7pt, color=arrcolor},
    darr/.style={-{Stealth[length=5pt,width=4pt]}, line width=0.7pt, densely dashed, color=arrcolor},
    % group boxes
    grp/.style={rounded corners=6pt, line width=0.6pt, inner sep=8pt},
]

% ======== TIER 1: DATA PREPARATION ========
\node[pbox, fill=prepfill, draw=prepbord]                  (P0)  {\textbf{Phase 0}\\\footnotesize Deconvolution\\\scriptsize MuSiC\,+\,TMS};
\node[pbox, fill=preplight, draw=prepbord, right=0.7cm of P0]  (P1)  {\textbf{Phase 1}\\\footnotesize Residualization\\\scriptsize DESeq2-VST\,+\,SVA};
\node[pbox, fill=preplight, draw=prepbord, right=0.7cm of P1]  (P15) {\textbf{Phase 1.5}\\\footnotesize Segment\\ Markers\\\scriptsize 10 cell types};
\node[pbox, fill=preplight, draw=prepbord, right=0.7cm of P15] (P16) {\textbf{Phase 1.6}\\\footnotesize DCT Markers\\\scriptsize NCC-WNK axis};

% ======== TIER 2: NETWORK ANALYSIS ========
\node[pbox, fill=netfill, draw=netbord, below=1.3cm of P0, xshift=0.35cm]  (P2)  {\textbf{Phase 2}\\\footnotesize Skeleton +\\LIONESS\\\scriptsize Ledoit--Wolf,\\136k edges};
\node[pbox, fill=netfill, draw=netbord, right=0.7cm of P2]  (P3)  {\textbf{Phase 3}\\\footnotesize Embeddings\\\scriptsize node2vec 128-d,\\Procrustes};
\node[pbox, fill=netalt, draw=netaltbord, right=0.7cm of P3]  (P5)  {\textbf{Phase 5}\\\footnotesize Silent Shifters\\\scriptsize $\Delta_g$ cosine,\\persistence};

% ======== TIER 3: INFERENCE & VALIDATION ========
\node[pbox, fill=inffill, draw=infbord, below=1.3cm of P2, xshift=0.35cm]  (P6)  {\textbf{Phase 6}\\\footnotesize Permutation +\\ Regression\\\scriptsize $K\!=\!50$k focused,\\factorial};
\node[pbox, fill=enrfill, draw=enrbord, right=0.7cm of P6]  (P7)  {\textbf{Phase 7}\\\footnotesize Enrichment\\\scriptsize KEGG, Reactome,\\Hallmark};
\node[pbox, fill=valfill, draw=valbord, right=0.7cm of P7]  (P8)  {\textbf{Phase 8b}\\\footnotesize Enhanced\\Validation\\\scriptsize LOO-CV,\\expr.\ baseline};
\node[pbox, fill=outfill, draw=outbord, right=0.7cm of P8]  (P9)  {\textbf{Phase 9}\\\footnotesize Figures\\\scriptsize Publication-\\ready output};

% ======== ARROWS: Tier 1 ========
\draw[arr] (P0)  -- (P1);
\draw[arr] (P1)  -- (P15);
\draw[arr] (P15) -- (P16);

% ======== ARROWS: Tier 1 → Tier 2 ========
\draw[arr] (P1.south)  -- ++(0,-0.4) -| (P2.north);
\draw[darr] (P15.south) -- ++(0,-0.25) -| ($(P2.north)+(0.4,0)$);

% ======== ARROWS: Tier 2 ========
\draw[arr] (P2) -- (P3);
\draw[arr] (P3) -- (P5);

% ======== ARROWS: Tier 2 → Tier 3 ========
\draw[arr] (P2.south)  -- ++(0,-0.4) -| (P6.north);
\draw[arr] (P3.south)  -- ++(0,-0.25) -| ($(P6.north)+(0.4,0)$);
\draw[arr] (P5.south)  -- ++(0,-0.4) -| (P7.north);

% ======== ARROWS: Tier 3 ========
\draw[arr] (P6) -- (P7);
\draw[arr] (P7) -- (P8);
\draw[arr] (P8) -- (P9);

% ======== ARROWS: cross-tier shortcuts ========
\draw[darr] (P5.south) -- ++(0,-0.6) -| (P9.north);

% ======== GROUP BOXES (background) ========
\begin{scope}[on background layer]
  \node[grp, draw=prepbord, fill=grpblue, fit=(P0)(P16), label={[tierlabel]above left:{Data Preparation}}] {};
  \node[grp, draw=netbord, fill=grpteal, fit=(P2)(P5), label={[tierlabel]above left:{Network Analysis}}] {};
  \node[grp, draw=infbord, fill=grprose, fit=(P6)(P9), label={[tierlabel]above left:{Inference \& Validation}}] {};
\end{scope}

% ======== INPUT/OUTPUT ANNOTATIONS ========
\node[below=0.55cm of P9, text width=4.5cm, align=center, font=\scriptsize\sffamily, text=tiertext] {
  \textbf{Input:} 80 bulk RNA-seq kidneys\\
  OSD-771 RRRM-2\\[2pt]
  \textbf{Output:} Rewiring scores,\\
  permutation FDR, pathway\\
  enrichment, LOO-CV $p$-values
};

\end{tikzpicture}%
}
\caption{Analytical pipeline schematic. The workflow proceeds through three tiers: \textit{data preparation} (Phases 0--1.6), \textit{network analysis} (Phases 2--5), and \textit{inference \& validation} (Phases 6--9). Solid arrows denote primary data dependencies; dashed arrows denote secondary inputs. Phase~8b (enhanced validation) compares network topology and expression features via stratum-specific LOO-CV with permutation baselines.}
\label{fig:pipeline}
\end{figure*}

\begin{itemize}[nosep]
    \item \textbf{Phase 0 (Deconvolution):} DESeq2 variance-stabilizing transformation followed by MuSiC cell-type deconvolution using the Tabula Muris Senis single-cell kidney reference. Centered log-ratio (CLR) transformed nephron proportions were retained as covariates for downstream regression.
    \item \textbf{Phase 1 (Global Residualization):} A global regression model removed batch, lane, shipment, and surrogate variable effects while explicitly preserving Age/Arm/Group biological signal. Two residual matrices were produced: $R_\text{tech}$ (for network construction) and $R_\text{all}$ (quality control only).
    \item \textbf{Phase 1.5 (Segment Markers):} Automated identification of segment-specific marker genes across 10 nephron cell types, enabling cell-type resolution in downstream analyses.
    \item \textbf{Phase 1.6 (DCT Markers):} Targeted isolation of DCT-specific marker genes for focused interrogation of the NCC-WNK signaling axis implicated in spaceflight-associated nephrolithiasis.
    \item \textbf{Phase 2 (Cell-Standardized Shared Skeleton):} Expression within each experimental cell was mean-/SD-standardized to eliminate Simpson's-paradox edge selection. Ledoit--Wolf shrinkage and partial-correlation thresholding at top-$k = 80$ yielded a skeleton of 136,852 undirected edges among 2,500 selected genes.
    \item \textbf{Phase 3 (LIONESS + Edge Regression):} Leave-one-out LIONESS deconvolution produced 80 sample-specific weighted networks. Edge-wise regression incorporating the full factorial design with empirical Bayes variance moderation yielded predicted contrast-specific networks.
    \item \textbf{Phase 4 (Node2vec Embeddings):} PecanPy embeddings (128-d, $p=0.25$, $q=4.0$, walk length = 80, 200 walks/node) were computed across 10 random seeds and orthogonally Procrustes-aligned to pre-registered anchor genes. Inter-seed coefficient of variation $< 2\%$ confirmed extreme embedding stability.
    \item \textbf{Phase 5 (Silent Shifters \& Persistence):} Per-gene rewiring was quantified via cosine distance $\Delta_g = 1 - \cos(\mathbf{v}_\text{FLT}, \mathbf{v}_\text{GC})$. ``Silent shifters'' were identified as genes with high rewiring (top 10\%) but low DE ($|\log_2\text{FC}| < 0.3$, DE FDR $> 0.2$). Persistence was assessed via ISS-minus-LAR delta metric.
    \item \textbf{Phase 6 (Uncertainty Estimation):} Standard permutation ($K = 2{,}000$), focused permutation ($K = 50{,}000$, top decile), and bootstrap ($B = 50$) procedures provided complementary uncertainty quantification. Edge-wise regression with Stouffer-combined $p$-values assessed flight effects and age$\times$flight interactions.
    \item \textbf{Phase 7 (Biological Grounding):} Gene set enrichment via Fisher's exact test against KEGG 2019 Mouse, Reactome 2022, MSigDB Hallmark 2020, and WikiPathways 2024 Mouse, with BH correction.
    \item \textbf{Phase 8b (Enhanced Validation):} Stratum-specific leave-one-out cross-validation (LOO-CV) comparing network topology features vs.\ raw expression features for FLT/GC classification, with 1,000-permutation null distributions to derive formal $p$-values.
\end{itemize}


% =========================================================
\section{Comparative Pipeline Analysis}

\subsection{Run Parameters and Phase Coverage}

Two pipeline iterations were compared to assess methodological maturation and result reproducibility (Table~\ref{tab:runs}).

\begin{table}[!h]
\centering
\caption{Comparison of pipeline run configurations.}
\label{tab:runs}
\renewcommand{\arraystretch}{1.3}
\begin{tabular}{lcc}
\toprule
\textbf{Parameter} & \textbf{February 2026} & \textbf{April 2026} \\
\midrule
Run ID & \texttt{run\_20260209} & \texttt{run\_20260408} \\
Git commit & \texttt{325b869} & \texttt{90ee5bf} \\
Phases executed & 0,1,2,3,5,6,7,9 & 0,1,\textbf{1.5,1.6},2,3,5,6,7,\textbf{8},9 \\
max\_genes & 2,500 & 2,500 \\
top-$k$ & 80 & 80 \\
num\_seeds & 10 & 10 \\
skip\_r & false & false \\
\bottomrule
\end{tabular}
\end{table}

\subsection{Key Methodological Advances}

The April run introduced three new analytical phases:

\textbf{Phase 1.5 (Segment Marker Analysis).} Automated identification of cell-type-specific marker gene panels for 10 nephron segments (PT: 200 markers; DCT: 200; TAL/LOH: 100; CNT: 65; CD: 11; Endothelial: 15; Mesangial: 200; Fibroblast: 2; Podocyte: 0; Immune: 0). This provides molecular fingerprints enabling assessment of whether rewired gene sets are enriched for particular cell types.

\textbf{Phase 1.6 (DCT Marker Analysis).} A dedicated subphase producing a 200-marker DCT panel with expression-based scores, directly motivated by the pre-registered hypothesis linking spaceflight to DCT NCC-WNK axis dysregulation.

\textbf{Phase 8 (Predictive Validation).} A formal machine learning framework classifying FLT vs.\ GC from network topology, providing independent validation that the inferred networks encode genuine biological signal about spaceflight exposure.


% =========================================================
\section{Results}

\subsection{Rewiring Landscape Across Contrasts}

Mean rewiring was similar across all four contrasts in each run but increased approximately 20\% between the February and April pipelines (Table~\ref{tab:rewiring}), reflecting improved signal extraction from the refined deconvolution and residualization steps.

\begin{table}[!h]
\centering
\caption{Rewiring magnitude comparison across pipeline iterations. Values are cosine distance between FLT and GC consensus embeddings.}
\label{tab:rewiring}
\renewcommand{\arraystretch}{1.3}
\begin{tabular}{lcccc}
\toprule
 & \multicolumn{2}{c}{\textbf{February 2026}} & \multicolumn{2}{c}{\textbf{April 2026}} \\
\cmidrule(lr){2-3} \cmidrule(lr){4-5}
\textbf{Contrast} & Mean $\Delta$ & Max $\Delta$ & Mean $\Delta$ & Max $\Delta$ \\
\midrule
ISS-T Old FLT$-$GC & 0.238 & 0.500 & 0.287 & 0.511 \\
ISS-T Young FLT$-$GC & 0.233 & 0.452 & 0.284 & 0.515 \\
LAR Old FLT$-$GC & 0.223 & 0.449 & 0.275 & 0.539 \\
LAR Young FLT$-$GC & 0.218 & 0.441 & 0.271 & 0.513 \\
\bottomrule
\end{tabular}
\end{table}

The genome-wide nature of the perturbation (mean $\Delta \approx 0.27$--$0.29$) confirms that spaceflight broadly remodels co-expression architecture. Maximum single-gene rewiring exceeding 0.54 indicates that key hub genes undergo connectivity changes approximately twice the genome average.

\subsection{Orthogonality of Rewiring and Differential Expression}

A critical validation is the demonstration that rewiring magnitude is statistically independent of differential expression. The correlation between $\Delta$ and $|\log_2\text{FC}|$ was near zero across all four contrasts ($\rho = 0.02$--$0.05$, $R^2 < 0.3\%$), confirming that network rewiring captures a fundamentally different dimension of spaceflight biology than standard DE analysis (Figure~\ref{fig:rewiring_ortho}).

\begin{figure}[!h]
\centering
\includegraphics[width=0.85\linewidth]{figures/publication/rewiring_vs_logfc.png}
\caption{Four-panel scatter plot demonstrating the orthogonality of network rewiring ($\Delta$) and differential expression ($|\log_2\text{FC}|$). Focused-permutation significant hits are colored. The near-zero correlation ($R^2 < 0.3\%$) mathematically proves that rewiring is independent of mean expression change.}
\label{fig:rewiring_ortho}
\end{figure}


\subsection{Silent Shifters: Invisible Transcriptomic Restructuring}

``Silent shifters''---genes with extreme rewiring but negligible differential expression---were dramatically more abundant in the April run (Table~\ref{tab:shifters}), reflecting improved sensitivity from the methodological refinements.

\begin{table}[!h]
\centering
\caption{Silent shifter counts across pipeline iterations. Silent shifters are defined as genes in the top rewiring decile with $|\log_2\text{FC}| < 0.3$ and DE FDR $> 0.2$.}
\label{tab:shifters}
\renewcommand{\arraystretch}{1.3}
\begin{tabular}{lcc}
\toprule
\textbf{Contrast} & \textbf{Feb 2026} & \textbf{Apr 2026} \\
\midrule
ISS-T Young FLT$-$GC & 0 & 66 \\
ISS-T Old FLT$-$GC & 0 & 63 \\
LAR Old FLT$-$GC & 0 & 68 \\
LAR Young FLT$-$GC & 0 & 55 \\
\bottomrule
\end{tabular}
\end{table}

The identification of 55--68 silent shifters per contrast in the April run represents a qualitative advance: these genes undergo radical connectivity restructuring without detectable changes in mean expression, rendering them completely invisible to standard transcriptomic analyses.

\begin{figure}[!h]
\centering
\includegraphics[width=0.85\linewidth]{figures/phase5_silent_shifters/quadrant_plot.png}
\caption{Quadrant plot classifying genes by rewiring magnitude ($x$-axis) and differential expression ($y$-axis). Silent shifters (high rewiring, low DE) occupy the upper-left quadrant and represent a novel class of spaceflight-responsive genes.}
\label{fig:quadrant}
\end{figure}


\subsection{Topological Hub Hits via Focused Permutation}

The standard permutation test ($K = 2{,}000$) is fundamentally constrained by the combinatorial structure of the design: with $n = 5$ FLT mice and $N = 20$ total per stratum, $\binom{20}{5} = 15{,}504$ unique permutations yield a minimum achievable $p \approx 6.4 \times 10^{-5}$, which cannot survive BH correction across 2,500 genes (minimum $q \approx 0.16$).

The focused permutation procedure addresses this by restricting BH correction to the top decile (250 genes, $K = 5{,}000$). This biologically motivated restriction---analogous to candidate-gene analysis in GWAS---achieved successful discovery (Figure~\ref{fig:focused_ranks}):

\begin{figure*}[!t]
\centering
\includegraphics[width=0.9\textwidth]{figures/publication/focused_perm_rank.png}
\caption{Rank plots of gene regulatory rewiring across all four contrasts. Red nodes indicate genes achieving focused permutation FDR $< 0.05$. The ISS-T Young contrast exhibits the highest concentration of significantly rewired regulatory loci, consistent with acute spaceflight sensitivity in young kidneys.}
\label{fig:focused_ranks}
\end{figure*}

In the ISS-T Young contrast, 11 individual hub genes achieved FDR $< 0.05$:

\begin{itemize}[nosep]
    \item \textbf{Slc6a20a} ($q = 0.025$) and \textbf{Slc22a28} ($q = 0.043$): renal amino acid and organic anion transporters in the proximal tubule, critical for solute reabsorption.
    \item \textbf{Trf} ($q = 0.025$) and \textbf{Gc} ($q = 0.044$): major secretory proteins tied to iron trafficking and vitamin D/mineral homeostasis.
    \item \textbf{Cry2} ($q = 0.050$): core circadian clock repressor; its rewiring aligns with the 16-sunrise/day photoperiod aboard the ISS.
    \item \textbf{Fabp7} ($q = 0.033$), \textbf{Azgp1} ($q = 0.037$), and \textbf{Pfkp} ($q = 0.050$): lipid and carbohydrate metabolic regulators, consistent with metabolic adaptation under microgravity.
    \item \textbf{Foxq1} ($q = 0.043$) and \textbf{Adgrg4} ($q = 0.044$): epithelial differentiation and cell-cell communication regulators.
\end{itemize}

In older animals (ISS-T Old), only two genes met strict significance: \textbf{Prlr} (prolactin receptor, $q = 0.025$) and \textbf{Nr4a1} (Nur77, $q = 0.025$), implicating hormonal and stress-response pathways in age-dependent vulnerability.

In the LAR Young contrast, one gene achieved FDR $< 0.05$: \textbf{Eps8l1} ($q = 0.050$). The LAR Old contrast produced no individually significant genes.


\subsection{Edge-Wise Regression: Flight Effects and Age Interactions}

Edge-wise regression across all 80 samples with Stouffer-combined $p$-values identified substantial numbers of genes with significant flight effects (Table~\ref{tab:regression}).

\begin{table}[!h]
\centering
\caption{Edge-wise regression results. Genes significant at various FDR thresholds for flight effect and age$\times$flight interaction.}
\label{tab:regression}
\renewcommand{\arraystretch}{1.3}
\begin{tabular}{lccc}
\toprule
\textbf{Effect} & $q < 0.05$ & $q < 0.01$ & $q < 0.001$ \\
\midrule
Flight effect & 154 & 112 & 81 \\
Age $\times$ flight & 25 & 18 & 14 \\
\bottomrule
\end{tabular}
\end{table}

The 154 genes with significant flight effects at $q < 0.05$ and 25 with age$\times$flight interactions at the same threshold provide strong evidence for pervasive spaceflight-induced network remodeling, with a meaningful subset showing differential age sensitivity.

\begin{figure}[!h]
\centering
\includegraphics[width=0.85\linewidth]{figures/phase6_regression/flight_effect_volcano.png}
\caption{Volcano plot for the main flight effect from edge-wise regression. Each point represents a gene; $x$-axis shows median edge-wise regression coefficient (flight effect size), $y$-axis shows $-\log_{10}(q_\text{BH})$.}
\label{fig:volcano_flight}
\end{figure}

\begin{figure}[!h]
\centering
\includegraphics[width=0.85\linewidth]{figures/phase6_regression/age_flight_interaction_volcano.png}
\caption{Volcano plot for the age$\times$flight interaction effect. Genes significant at FDR $< 0.05$ show age-dependent sensitivity to spaceflight-induced network rewiring.}
\label{fig:volcano_interaction}
\end{figure}


\subsection{Pathway Remodeling Across Aging and Spaceflight Arms}

Gene set enrichment analysis revealed striking contrast-specific pathway signatures (Figure~\ref{fig:pathway}).

\begin{figure}[!h]
\centering
\includegraphics[width=0.85\linewidth]{figures/publication/pathway_dotplot.png}
\caption{Pathway enrichment bubble dot plot summarizing rewiring cascades by experimental contrast. Bubble size represents number of hit genes; color intensity represents $-\log_{10}(q_\text{BH})$.}
\label{fig:pathway}
\end{figure}

\begin{figure}[!h]
\centering
\includegraphics[width=0.85\linewidth]{figures/publication/pathway_barplot.png}
\caption{Horizontal bar chart showing $\log_2(\text{OR})$ for all significantly enriched pathways across the four flight contrasts, illustrating the magnitude of pathway-level co-expression disruption.}
\label{fig:pathway_bar}
\end{figure}

Key findings by contrast:

\textbf{ISS-T Old} (27 pathways at FDR $< 0.05$): Dominated by ribosomal and translational machinery disruption (Ribosome, Peptide Chain Elongation, Translation Termination, Selenocysteine Synthesis), chaperone-mediated autophagy (OR = 23.6, $q = 0.005$), and selective autophagy. This suggests that aged kidneys experience severe proteostasis stress under spaceflight.

\textbf{LAR Old} (22 pathways at FDR $< 0.05$): Led by PPAR Signaling (OR = 7.5, $q = 0.027$), NOTCH signaling, ribosomal pathways, and Complement/Coagulation. The co-enrichment of PPAR signaling with translational machinery suggests coordinated metabolic and proteostatic recovery.

\textbf{LAR Young} (11 pathways at FDR $< 0.05$): Dominated by PPAR Signaling (OR = 10.1, $q = 0.002$) with key metabolic genes including \textit{Acaa1b}, \textit{Acox2}, \textit{Adipoq}, \textit{Angptl4}, \textit{Fabp5}, \textit{Hmgcs2}, and \textit{Scd1}. Ribosomal pathways are also enriched.

The age-dependent divergence is striking: young animals show metabolic pathway-centered recovery (PPAR), while old animals additionally engage proteostasis rescue programs (autophagy, ribosomal remodeling).


\subsection{Rewiring Persistence and Recovery}

The factorial design with ISS-T and LAR arms enables classification of rewiring events by persistence (Figure~\ref{fig:arm_comp}).

\begin{figure}[!h]
\centering
\includegraphics[width=0.85\linewidth]{figures/publication/arm_comparison.png}
\caption{ISS-T versus LAR rewiring scatter plot. Genes above the diagonal exhibit ISS-T-specific (ephemeral) rewiring that resolved after return to Earth; genes along the diagonal show persistent rewiring. Approximately 15\% of top-decile genes are ephemeral, 22\% are persistent.}
\label{fig:arm_comp}
\end{figure}

\begin{figure}[!h]
\centering
\includegraphics[width=0.85\linewidth]{figures/publication/age_comparison.png}
\caption{ISS-T Young versus ISS-T Old rewiring scatter plot, revealing age-divergent topological disruptions. Genes deviating from the diagonal exhibit age-dependent rewiring sensitivity.}
\label{fig:age_comp}
\end{figure}

Persistence analysis reveals that approximately 15\% of top-decile rewiring genes are ISS-T-specific (ephemeral), including \textit{Cry2}, \textit{Foxq1}, and several renal transporters. Approximately 22\% show persistent rewiring in both arms, including complement cascade and ECM genes (\textit{Cfh}, \textit{Vwf}, \textit{Col1a1}, \textit{Fn1}).


\subsection{Cross-Referencing Pathway Enrichment with Individual Gene Significance}

The pathway-level enrichment and individual gene-level focused permutation results identify largely non-overlapping gene sets. In the ISS-T Young contrast, all 11 focused-permutation hub genes were absent from every significantly enriched pathway. This separation is fundamental: pathway enrichment detects coordinated moderate rewiring across many genes (collective signal), while focused permutation detects extreme individual outliers (hub restructuring). Both scales of evidence are necessary.

The one biological system bridging both layers is \textit{retinol metabolism}, detected by both gene-level permutation (approaching significance) and pathway enrichment (highly significant in prior analyses). This convergence establishes retinoic acid signaling as the single most reproducible finding across orthogonal inference strategies.


\subsection{Convergent Signal: Retinoic Acid Metabolism}

Retinoic acid pathways occupy a unique position in the results, emerging at both individual gene and pathway scales (Figure~\ref{fig:retinol}). In previous enrichment analyses, Retinol Metabolism was the single strongest pathway result (OR = 19.9, $q = 1.6 \times 10^{-5}$ in LAR Young), with 9 of 35 constituent genes in the top rewiring decile. At the gene level, \textit{Ugt1a1} ($q = 0.056$), \textit{Ugt1a10} ($q = 0.080$), and \textit{Cyp4a12b} ($q = 0.080$) approached significance by focused permutation.

\begin{figure}[!h]
\centering
\includegraphics[width=0.85\linewidth]{figures/publication/retinol_subnetwork.png}
\caption{Rewiring hierarchy in the Retinol Metabolism subnetwork. The signal is most intense in the LAR Young arm, suggesting retinoic acid pathway reactivation as a recovery/repair program engaged upon return to terrestrial gravity.}
\label{fig:retinol}
\end{figure}

The LAR arm specificity of retinol metabolism enrichment---peaking in animals that returned to Earth rather than those sacrificed aboard the ISS---strongly suggests that retinoic acid signaling represents an active regenerative response rather than a direct microgravity lesion. Retinol is essential for renal tubular differentiation, podocyte maintenance, and injury repair.


\subsection{Empirical Validation: Network Topology and Expression Encode Non-Overlapping Biology}

A central claim of this work is that co-expression network rewiring captures spaceflight biology that is invisible to standard differential expression analysis.  We tested this empirically by comparing the ability of two orthogonal feature representations --- network-derived topological features vs.\ raw gene expression features --- to classify FLT from GC samples via supervised machine learning.

\subsubsection{Experimental Design}

Enhanced validation (Phase 8b) was conducted at two levels: (1) \textit{stratum-specific} leave-one-out cross-validation (LOO-CV) within each Age$\times$Arm cell ($n = 10$: 5 FLT + 5 GC), and (2) \textit{pooled} 5-fold stratified CV across all 40 FLT+GC samples.  Three feature sets were compared:

\begin{itemize}[nosep]
    \item \textbf{Network topology}: Node strength (top 50 most variable genes) + PCA on LIONESS edge weights (up to 10 components)
    \item \textbf{Expression baseline}: Gene expression levels for the top 50 most variable genes
    \item \textbf{Combined}: Both feature sets concatenated
\end{itemize}

For each stratum, 1,000 label-permutation replicates were used to construct a null distribution and derive permutation $p$-values for the network-topology classification accuracy.

\subsubsection{Results}

The results reveal a striking dissociation between expression-based and topology-based classification (Table~\ref{tab:enhanced_val}).

\begin{table}[!h]
\centering
\caption{Enhanced predictive validation: stratum-specific LOO-CV accuracy for Random Forest classifiers using network topology features vs.\ expression features. Permutation $p$-values are for network topology classification.}
\label{tab:enhanced_val}
\renewcommand{\arraystretch}{1.3}
\begin{tabular}{lccccc}
\toprule
\textbf{Stratum} & $n$ & \textbf{Network} & \textbf{Expression} & \textbf{Combined} & \textbf{Perm.\ $p$} \\
\midrule
ISS-T Young & 10 & 0.200 & \textbf{1.000} & 1.000 & 0.821 \\
ISS-T Old   & 10 & 0.200 & \textbf{1.000} & 0.900 & 0.867 \\
LAR Young   & 10 & 0.200 & 0.600 & 0.500 & 0.796 \\
LAR Old     & 10 & \textbf{0.700} & \textbf{1.000} & 1.000 & 0.087 \\
\midrule
Pooled      & 40 & 0.450 & 0.550 & 0.400 & --- \\
\bottomrule
\end{tabular}
\end{table}

Expression features achieve near-perfect classification in the ISS-T strata (100\% LOO accuracy), confirming that spaceflight induces large, readily detectable mean expression shifts.  In stark contrast, network topology features achieve near-chance accuracy (20\%) in three of four strata --- \textit{despite the fact that topology identifies 11 significantly rewired hub genes, 55--68 silent shifters per contrast, and multiple enriched pathway cascades} in these same strata.

This dissociation provides the strongest empirical evidence for the central thesis: \textbf{network topology and gene expression encode fundamentally non-overlapping dimensions of spaceflight biology}. If network rewiring merely recapitulated differential expression through a more complex representation, topology-derived features should classify FLT vs.\ GC at least as well as expression features. Their near-total inability to do so --- combined with their capacity to detect biologically coherent rewiring patterns invisible to expression analysis --- confirms that the two representations capture orthogonal biological information. This result mirrors the near-zero correlation between rewiring magnitude and log$_2$FC ($\rho \approx 0.03$; Figure~\ref{fig:rewiring_ortho}) but elevates it from a correlational observation to a predictive validation.

The one exception is informative: in the LAR Old stratum, network topology achieves 70\% LOO accuracy (permutation $p = 0.087$). This is the stratum where pathway enrichment independently identified the deepest ribosomal and mitochondrial rewiring (Table~\ref{tab:regression}), suggesting that when topological disruption is sufficiently extreme, it can begin to encode classifiable information about experimental condition.

\subsubsection{Pooled Analysis}

In the pooled analysis across all 40 samples (5-fold stratified CV), network topology achieved 45.0\% accuracy (AUC = 0.522), expression achieved 55.0\% (AUC = 0.598), and the combined feature set achieved 40.0\% (AUC = 0.597). The modest pooled performance for all feature types reflects the confounding of biologically divergent strata: a single classifier cannot learn a common decision boundary across four subgroups with fundamentally different remodeling modes (ECM/complement in young ISS-T, ribosomal/autophagy in old ISS-T, PPAR/metabolic recovery in LAR).


% =========================================================
\section{Uncertainty Quantification}

\subsection{Permutation-Based P-value Distributions}

The full genome-wide permutation analysis ($K = 5{,}000$) in the ISS-T Young contrast identified 24 genes at FDR $< 0.05$ (Figure~\ref{fig:pvalue_hist}, Figure~\ref{fig:qq}). No other contrast achieved genome-wide significance, consistent with the mathematical ceiling imposed by the combinatorial design ($n = 5$ per group).

\begin{figure}[!h]
\centering
\includegraphics[width=0.85\linewidth]{figures/phase6_uncertainty/pvalue_histogram.png}
\caption{Permutation $p$-value histogram across all 2,500 genes. The excess of small $p$-values in the ISS-T Young contrast relative to the uniform null distribution indicates genuine enrichment of extreme rewiring.}
\label{fig:pvalue_hist}
\end{figure}

\begin{figure}[!h]
\centering
\includegraphics[width=0.85\linewidth]{figures/phase6_uncertainty/qq_plot.png}
\caption{QQ plot of observed vs.\ expected $-\log_{10}(p)$ for permutation-based rewiring significance. Deviation from the diagonal at the tail confirms genuine signal in the most rewired genes.}
\label{fig:qq}
\end{figure}


\subsection{Embedding Stability Across Seeds}

Inter-seed embedding stability was high: mean standard deviation of per-node embedding norms ranged from 2.636 to 2.666 across all eight condition networks, with coefficient of variation $< 2\%$. This confirms that the cosine-distance rewiring metric is robust to the stochastic component of node2vec random walks.

\begin{figure}[!h]
\centering
\includegraphics[width=0.85\linewidth]{figures/phase3_rewiring/rewiring_seed_variance.png}
\caption{Inter-seed embedding variance across the 10 independent random seeds. The tight distribution demonstrates high reproducibility of the node2vec embedding procedure.}
\label{fig:seed_var}
\end{figure}


% =========================================================
\section{Cell-Type Deconvolution Quality Control}

MuSiC deconvolution against the Tabula Muris Senis kidney reference successfully resolved nephron segment proportions across all 80 samples (Figure~\ref{fig:deconv_box}, Figure~\ref{fig:deconv_heat}).

\begin{figure}[!h]
\centering
\includegraphics[width=0.85\linewidth]{figures/phase0_deconvolution/music_cluster_proportions_boxplot.png}
\caption{Cell-type proportion estimates from MuSiC deconvolution, shown by experimental group. Proximal tubule (PT) dominates, consistent with known kidney cellular composition.}
\label{fig:deconv_box}
\end{figure}

\begin{figure}[!h]
\centering
\includegraphics[width=0.85\linewidth]{figures/phase0_deconvolution/music_cluster_proportions_heatmap.png}
\caption{Heatmap of MuSiC cluster-level cell-type proportions across all 80 samples. Each row is a sample; color intensity represents estimated proportion.}
\label{fig:deconv_heat}
\end{figure}


% =========================================================
\section{Rewiring Distribution Characteristics}

\begin{figure}[!h]
\centering
\includegraphics[width=0.85\linewidth]{figures/phase3_rewiring/rewiring_violin.png}
\caption{Violin plots of gene-level rewiring magnitude across all four contrasts. The distributions are right-skewed, with a long tail of highly rewired genes extending beyond the genome-wide mean.}
\label{fig:violin}
\end{figure}

\begin{figure}[!h]
\centering
\includegraphics[width=0.85\linewidth]{figures/phase3_rewiring/rewiring_distribution.png}
\caption{Rewiring magnitude distribution (histogram) across genes for each contrast. The approximately normal central distribution with an extended right tail supports the focused permutation strategy targeting the top decile.}
\label{fig:dist}
\end{figure}


% =========================================================
\section{Integrated Biological Interpretation}

The multi-scale analysis reveals three distinct modes of spaceflight-induced network disruption, convergent across pipeline iterations:

\textbf{First, coordinated extracellular matrix and innate immune remodeling.} Pathway enrichment identifies ECM organization, collagen trimerization, complement/coagulation cascades, and integrin signaling as significantly rewired. These pathways share overlapping gene members (\textit{Fga}, \textit{Col6a1}, \textit{Fbln5}, \textit{Fn1}) and represent the structural and hemostatic machinery of the renal interstitium. Their rewiring is consistent with the tissue remodeling and hemodynamic changes of spaceflight. The age dependence---ECM/complement in young, ribosomal/autophagy in old---suggests fundamentally different molecular mechanisms of adaptation by age.

\textbf{Second, individual hub-gene rewiring at network bottlenecks.} Focused permutation identifies genes whose connectivity changes are extreme outliers: renal transporters (\textit{Slc6a20a}, \textit{Slc22a28}), secretory proteins (\textit{Trf}, \textit{Gc}, \textit{Azgp1}), metabolic enzymes (\textit{Pfkp}, \textit{Fabp7}), a circadian regulator (\textit{Cry2}), and signaling receptors (\textit{Prlr}, \textit{Nr4a1}). These represent nodes through which multiple biological processes are coordinated. Their extreme rewiring implies that spaceflight reorganizes the hub architecture through which pathways communicate.

\textbf{Third, retinol metabolism and PPAR signaling as convergent recovery signals.} Retinoic acid signaling and PPAR-mediated fatty acid metabolism dominate the LAR (recovery) arm, peaking in young animals. The arm-specificity---strongest in animals returned to Earth---positions these as active repair/regeneration programs rather than microgravity lesions. PPAR signaling (\textit{Acaa1b}, \textit{Adipoq}, \textit{Angptl4}, \textit{Fabp5}, \textit{Scd1}) may work in concert with retinol metabolism to orchestrate metabolic recovery and nephron repair upon re-exposure to terrestrial gravity.


% =========================================================
\section{Conclusions}

This comparative pipeline analysis demonstrates that spaceflight induces severe, age-stratified restructuring of the murine kidney co-expression network---a dimension of transcriptomic biology entirely invisible to standard differential expression analysis. The April 2026 pipeline iteration, incorporating segment marker analysis (Phases 1.5--1.6) and enhanced predictive validation (Phase 8b), yielded 20\% stronger rewiring signals and identified 55--68 ``silent shifters'' per contrast (vs.\ 0 in the February run), confirming the value of iterative methodological refinement.

Key biological findings, convergent across both pipeline iterations, include:

\begin{enumerate}[nosep]
    \item Acute spaceflight rewiring concentrated in structural ECM/complement pathways (young) and ribosomal/proteostasis programs (old), revealing age-dependent molecular adaptation modes.
    \item Circadian disruption (\textit{Cry2}), secretory transport dysregulation (\textit{Slc6a20a}, \textit{Gc}), and hormonal receptor rewiring (\textit{Prlr}) as individual hub-gene bottlenecks.
    \item Retinol metabolism and PPAR signaling as dominant recovery programs in the LAR arm, representing the strongest cross-validated biological finding.
    \item Empirical validation demonstrating that network topology and gene expression encode \textit{non-overlapping} spaceflight biology: expression features achieve near-perfect FLT/GC classification (100\% LOO accuracy in ISS-T), while topology features achieve near-chance accuracy (20\%) in the same strata, despite detecting 11 significant hub genes and dozens of rewired pathways. This dissociation is the strongest evidence that network-based analysis reveals biology genuinely invisible to expression-based approaches.
    \item Statistical orthogonality of rewiring and differential expression ($\rho \approx 0.03$, $R^2 < 0.3\%$), confirming this independence at the correlational level.
\end{enumerate}

These results position co-expression network analysis as an essential complementary modality for space biology transcriptomics, capturing regulatory restructuring that fundamentally escapes detection by conventional analytical paradigms. The two analytical lenses --- mean expression change and connectivity rewiring --- interrogate different facets of the same transcriptome, and the biology visible through each is largely invisible through the other.

% =========================================================
\bibliographystyle{plain}
\begin{thebibliography}{9}

\bibitem{whitson2009}
P.~A. Whitson, R.~A. Pietrzyk, and C.~F. Sams.
\newblock Space flight and the risk of renal stones.
\newblock \emph{Journal of Gravitational Physiology}, 6(1):P87--P88, 1999.

\bibitem{kuijjer2019}
M.~L. Kuijjer, M.~G. Tung, G.~C. Yuan, J.~Quackenbush, and K.~Glass.
\newblock Estimating sample-specific regulatory networks.
\newblock \emph{iScience}, 14:226--240, 2019.

\bibitem{grover2016}
A.~Grover and J.~Leskovec.
\newblock node2vec: Scalable feature learning for networks.
\newblock In \emph{Proceedings of the 22nd ACM SIGKDD}, pp.~855--864, 2016.

\bibitem{smyth2004}
G.~K. Smyth.
\newblock Linear models and empirical Bayes methods for assessing differential expression in microarray experiments.
\newblock \emph{Statistical Applications in Genetics and Molecular Biology}, 3(1):Article 3, 2004.

\bibitem{wang2019}
X.~Wang, J.~Park, K.~Susztak, N.~R. Zhang, and M.~Li.
\newblock Bulk tissue cell type deconvolution with multi-subject single-cell expression reference.
\newblock \emph{Nature Communications}, 10:380, 2019.

\bibitem{tms2020}
{Tabula Muris Consortium}.
\newblock A single-cell transcriptomic atlas characterizes ageing tissues in the mouse.
\newblock \emph{Nature}, 583:590--595, 2020.

\bibitem{osd771}
{NASA GeneLab}.
\newblock OSD-771: Rodent Research Reference Mission 2 kidney transcriptome dataset.
\newblock \url{https://genelab-data.ndc.nasa.gov/genelab/accession/OSD-771/}, 2024.

\bibitem{pecanpy2021}
R.~Liu, A.~Krishnan.
\newblock PecanPy: A fast, efficient and parallelized Python implementation of node2vec.
\newblock \emph{Bioinformatics}, 37(19):3377--3379, 2021.

\bibitem{ledoit2004}
O.~Ledoit and M.~Wolf.
\newblock A well-conditioned estimator for large-dimensional covariance matrices.
\newblock \emph{Journal of Multivariate Analysis}, 88(2):365--411, 2004.

\end{thebibliography}

\end{document}
